\documentclass[10pt,a4paper]{amsart}
\usepackage[margin=19mm]{geometry}
\usepackage{amsmath,amssymb,mathtools}
\usepackage[scr=rsfs,cal=euler]{mathalfa}
\usepackage{xfrac}
\usepackage[unicode,hidelinks,bookmarks=false]{hyperref}
\newcommand{\ie}{i.e.\ }
\newcommand{\cf}{cf.\ }
\newcommand{\resp}{respectively\ }
\newcommand{\Subsection}{\subsection}
\providecommand{\nobreakdash}{\nobreak\hbox{-}}
\setlength{\emergencystretch}{2em}
\multlinegap0pt
\usepackage{bm}
\usepackage{bbm}
\usepackage{dsfont}
\usepackage{xspace}
\usepackage{cancel}
\usepackage{mathtools}
\usepackage{amsthm}
\makeatletter
\@ifundefined{thm}{\newtheorem{thm}{Theorem}}{}
\makeatother
\title[Explicit Factorization of $X^n-1$ over $\mathbb{Z}\_{p^e}$ vi]{Explicit Factorization of $X^n-1$ over $\mathbb{Z}\_{p^e}$ via Cofactor-Free Single-Seed Hensel Lifting}
\author{Yongchao Wang, Yang Ding, Jiansheng Yang, Zhiqiu Huang}
\date{Source: arXiv:2606.20633v1 (CC BY 4.0)}
\begin{document}
\maketitle

\noindent\emph{Source and licence.} Explicit Factorization of $X^n-1$ over $\mathbb{Z}_{p^e}$ via Cofactor-Free Single-Seed Hensel Lifting. Version arXiv:2606.20633v1 (CC BY 4.0), distributed under the Creative Commons Attribution 4.0 International licence, as recorded in the arXiv licence metadata for this exact version and shown on the abstract page https://arxiv.org/abs/2606.20633v1. Full rights notice: licenses/arxiv-2606.20633-CC-BY-4.0.txt.\par\smallskip

\noindent\textit{Editorial scope.} Counted unit: the Thm whose source label is \texttt{thm:cofactor\_free} in the article (source0.tex lines 420--447, source label \texttt{thm:cofactor\_free}), delivered together with the article's own complete original proof of it (source0.tex lines 449--483, one \texttt{proof} environment of 35 lines). No mathematics is written, repaired or re-derived: statement and proof are the article's own text line for line. Required context retained uncounted: none. Every definition, display and label of those lines is retained unchanged. Omitted from this package and not redistributed: 1--419, 448--448, 484--605. Nothing inside a retained excerpt is omitted or altered: every retained body is delivered exactly as the article's lines are written, so no omission receipt of any kind accompanies this package. Citations: the retained excerpts carry no citation command at all, so no bibliography entry of the article is transcribed and no reference is redistributed. Reference adaptation: none was needed and none was made; every reference inside the delivered unit resolves inside it, so no unresolved reference or citation is produced. Printed slips kept literal: no printed word, symbol, hypothesis or number of the counted statement or of its proof was altered. Layout and class: the article's own class is \texttt{article} with packages including arxiv, amsmath, amssymb, amsfonts, amsthm, algorithm2e, graphicx, float, xcolor, enumitem, url, tikz; the wrapper is the lane's pinned \texttt{amsart} editorial wrapper at 10pt on A4, which restates the theorem-like environments the retained text needs and adds the article's own macro definitions that the retained text needs (none), with extra wrapper packages (bm, bbm, dsfont, xspace, cancel, mathtools). The delivered numbering is the unit's own. Assets: no retained span contains a figure environment, a graphics-inclusion command, a PDF-page inclusion, a table or a TikZ picture; no image file is copied into this package, the package declares no asset dependency and no image file is redistributed with it. Licence: version arXiv:2606.20633v1 of this article is distributed under the Creative Commons Attribution 4.0 International licence, as recorded in the arXiv licence metadata for this exact version and shown on the abstract page https://arxiv.org/abs/2606.20633v1. A full rights notice is delivered at licenses/arxiv-2606.20633-CC-BY-4.0.txt. Characters: every retained excerpt is pure ASCII, so the delivered engine, XeLaTeX with its default Latin Modern text font, reports no missing character.
\begin{thm}[Cofactor-Free Hensel Lift] \label{thm:cofactor_free}
Let $G_1(X) \in \mathbb{F}_p[X]$ be a monic irreducible factor of $X^n - 1$ of degree $m$, and let $H(X) = (X^n - 1)/G_1(X) \in \mathbb{F}_p[X]$ be the complementary cofactor. Define the cached inverse
\begin{equation}
    C(X) \;=\; \bigl[H(X)\bigr]^{-1} \bmod G_1(X) \quad \text{in } \mathbb{F}_p[X]/\langle G_1(X) \rangle
\end{equation}
which exists since $\gcd(G_1, H) = 1$ in $\mathbb{F}_p[X]$. Suppose $G_h(X) \in \mathbb{Z}_{p^h}[X]$ is a monic lift satisfying $X^n - 1 \equiv 0 \pmod{G_h(X), \, p^h}$. Then:

\begin{enumerate}
    \item[\textup{(1)}] The error polynomial at layer $h$,
    \begin{equation}
        E_h(X) \;\equiv\; \frac{X^n \bmod G_h(X) \;-\; 1}{p^h} \pmod{p}
    \end{equation}
    depends on the current approximation $G_h(X)$ and thus \emph{changes at every layer}.
    
    \item[\textup{(2)}] The next-layer lift is obtained by computing the correction in $\mathbb{F}_p[X]$:
    \begin{equation} \label{eq:cofactor_free_correction}
        \Delta G_h(X) \;=\; E_h(X) \cdot C(X) \bmod G_1(X)
    \end{equation}
    and updating the approximation in $\mathbb{Z}_{p^{h+1}}[X]$:
    \begin{equation} \label{eq:cofactor_free_update}
        G_{h+1}(X) \;=\; G_h(X) \;+\; p^h \cdot \Delta G_h(X) \bmod p^{h+1}
    \end{equation}
    which satisfies $X^n - 1 \equiv 0 \pmod{G_{h+1}(X), \, p^{h+1}}$.
    
    \item[\textup{(3)}] The cached inverse $C(X)$ is independent of the precision depth $h$: it is computed once over $\mathbb{F}_p$ and reused identically at every layer $h = 1, 2, \dots, e-1$.
\end{enumerate}
Consequently, after the one-time precomputation of $C(X)$, each lifting iteration requires only a polynomial multiplication and modular reduction in $\mathbb{F}_p[X]/\langle G_1(X) \rangle$, costing $O(m^2)$ per layer---independent of both the global degree $n$ and the precision depth $e$.
\end{thm}

\begin{proof}
Let $G_1(X)$ be a monic irreducible factor of $X^n-1$ over $\mathbb{F}_p$, with cofactor $H(X) = (X^n-1)/G_1(X) \in \mathbb{F}_p[X]$. Suppose that after $h-1$ lifting steps we have obtained a monic polynomial $G_h(X) \in \mathbb{Z}_{p^h}[X]$ satisfying $G_h(X) \equiv G_1(X) \pmod p$ and a quotient $Q_h(X)$ such that $X^n - 1 \equiv Q_h(X) G_h(X) \pmod{p^h}$. We prove that the correction $\Delta G(X)$ required to extend $G_h(X)$ to $\mathbb{Z}_{p^{h+1}}[X]$ is determined solely by objects in $\mathbb{F}_p$.

Performing polynomial division over $\mathbb{Z}_{p^{h+1}}[X]$:
\begin{equation}
    X^n - 1 = Q_h(X) G_h(X) + p^h E(X)
\end{equation}
where the remainder $R(X) = p^h E(X)$ has all coefficients divisible by $p^h$, and $E(X) \equiv R(X)/p^h \pmod p$ is the reduced error polynomial in $\mathbb{F}_p[X]$.

Define the target lifted factors as $G_{h+1}(X) = G_h(X) + p^h \Delta G(X)$ and $Q_{h+1}(X) = Q_h(X) + p^h \Delta Q(X)$. The required congruence is:
\begin{equation} \label{eq:lift_target}
    X^n - 1 \equiv \big(G_h(X) + p^h \Delta G(X)\big)\big(Q_h(X) + p^h \Delta Q(X)\big) \pmod{p^{h+1}}
\end{equation}
Expanding \eqref{eq:lift_target} and observing that the cross term $p^{2h} \Delta G \Delta Q$ vanishes modulo $p^{h+1}$ (since $2h \geq h+1$ for $h \geq 1$):
\begin{equation} \label{eq:lift_expanded}
    X^n - 1 \equiv Q_h(X) G_h(X) + p^h \big(Q_h(X) \Delta G(X) + G_h(X) \Delta Q(X)\big) \pmod{p^{h+1}}
\end{equation}
Substituting $X^n - 1 = Q_h(X) G_h(X) + p^h E(X)$ and cancelling $Q_h G_h$ yields:
\begin{equation}
    p^h E(X) \equiv p^h \big(Q_h(X) \Delta G(X) + G_h(X) \Delta Q(X)\big) \pmod{p^{h+1}}
\end{equation}
Dividing by $p^h$ projects the constraint into $\mathbb{F}_p[X]$. Since $Q_h(X) \equiv H(X) \pmod p$ and $G_h(X) \equiv G_1(X) \pmod p$:
\begin{equation} \label{eq:projected_error}
    E(X) \equiv H(X) \Delta G(X) + G_1(X) \Delta Q(X) \pmod p
\end{equation}
Evaluating \eqref{eq:projected_error} modulo the base seed ideal $G_1(X)$ annihilates the second term:
\begin{equation}
    E(X) \equiv H(X) \Delta G(X) \pmod{G_1(X)}
\end{equation}
Because $\gcd(G_1(X), H(X)) = 1$ in $\mathbb{F}_p[X]$, the cofactor $H(X)$ is invertible in the quotient ring $\mathbb{F}_p[X]/\langle G_1(X) \rangle$, yielding the unique solution:
\begin{equation}
    \Delta G(X) \equiv E(X) \cdot [H(X)]^{-1} \pmod{G_1(X)}
\end{equation}
This derivation holds identically for every layer $h \geq 1$, completing the induction.
\end{proof}

\end{document}
